\chapter{Typed Hypergraphs and Content Atoms}
\label{ch:typed-hypergraphs}

This chapter opens Part~II by introducing the algebraic primitives of the
Polyxan hyperstructure.  In analogy with how Chapter~\ref{ch:semantic-manifold}
constructed the semantic manifold $(\mathcal{M},g)$ and its tensorial
inhabitants, we now construct the discrete objects—\emph{typed hypergraphs} and
\emph{content atoms}—that constitute the categorical backbone of Polyxan media.

Typed hypergraphs provide the algebraic substrate on which semantic curvature,
media quines, Polycompiler rewrites, and eventually RSVP--Polyxan coupling
(Part~III) are defined.  The chapter culminates in the first algebraic
characterization of media quines as fixed points of the free-algebra monad
generated by content-atom insertion.

% ============================================================
\section{Content Atoms: The Polyxan Primitives}
% ============================================================

A \emph{content atom} is the discrete analogue of an RSVP field at a point.
It is the primitive semantic unit of the Polyxan theory.

\begin{definition}[Content Atom]
A content atom is a quadruple
\[
a = (\tau,\, \pi,\, \theta,\, w)
\]
consisting of:
\begin{enumerate}
\item a \emph{type} $\tau \in \mathbf{Type}$,
\item a \emph{payload} $\pi$ (any structured datum),
\item a \emph{tag} $\theta \in \mathbf{Tag}$ refining $\tau$,
\item a \emph{weight} $w \in \mathbb{R}_{\ge 0}$.
\end{enumerate}
\]
\end{definition}

Types classify atoms into semantically coherent families; tags refine them;
payloads store content; weights measure semantic salience or energetic mass.
This mirrors the RSVP decomposition $(\Phi,\mathbf{v},S)$, with $\theta$
playing the rôle of local entropy label.

\subsection{Typed Incidence Structure}

Content atoms will be arranged into hypergraphs whose vertices carry types and
tags.  Let $V$ denote the set of vertices, each equipped with a content atom
$a(v)$.  Hyperedges will enforce semantic relations and compatibility
conditions encoded in the type system.

% ============================================================
\section{Typed Hypergraphs: Definition and Structure}
% ============================================================

Polyxan hypergraphs generalize directed, labeled hypergraphs by adding
type, tag, and weight structure at both vertices and edges.

\begin{definition}[Typed Hypergraph]
A \emph{typed hypergraph} is a tuple
\[
\mathcal{G} = (V,E,\mathrm{src},\mathrm{tgt}, \mathrm{type},\mathrm{tag}, w)
\]
where
\begin{itemize}
\item $V$ is a set of vertices,
\item $E$ is a set of hyperedges,
\item $\mathrm{src}(e),\mathrm{tgt}(e) \subseteq V^{k_e}$ give the arity-$k_e$
  source and target signatures of $e$,
\item $\mathrm{type}: V\cup E \to \mathbf{Type}$ assigns semantic types,
\item $\mathrm{tag}: V\cup E \to \mathbf{Tag}$ refines types,
\item $w: V\cup E \to \mathbb{R}_{\ge 0}$ assigns semantic weights.
\end{itemize}
\end{definition}

The type assignment must satisfy the compatibility condition
\[
\mathrm{type}(e) \in \mathbf{Type}\big(\mathrm{src}(e),\mathrm{tgt}(e)\big),
\]
i.e.~edge types depend contravariantly and covariantly on the types of their
incident vertices.

Typed hypergraphs form the primary objects of Polyxan algebra.  They will
admit curvature, sheaf structure, spectral invariants, and (in Part~III)
semantic embeddings into $(\mathcal{M},g)$.

% ============================================================
\section{The Category TypHyper}
% ============================================================

\begin{definition}[Morphisms of Typed Hypergraphs]
A morphism $f : \mathcal{G} \to \mathcal{H}$ consists of maps
\[
f_V : V \to V_{\!H},
\qquad
f_E : E \to E_{\!H}
\]
preserving:
\begin{itemize}
\item the incidence structure:
  \[
    f_V(\mathrm{src}(e)) = \mathrm{src}(f_E(e)),\qquad
    f_V(\mathrm{tgt}(e)) = \mathrm{tgt}(f_E(e)),
  \]
\item types:
  \[
  \mathrm{type}_{\mathcal{G}}(x) \le_{\mathbf{Type}}
    \mathrm{type}_{\mathcal{H}}(f(x)),
  \]
\item tags and weights:
  \[
  \theta_{\mathcal{G}}(x) \le_{\mathbf{Tag}} \theta_{\mathcal{H}}(f(x)),
  \qquad
  w_{\mathcal{G}}(x) = w_{\mathcal{H}}(f(x)).
  \]
\end{itemize}
\end{definition}

\begin{proposition}
Typed hypergraphs and their morphisms form a category, denoted
\[
\mathbf{TypHyper}.
\]
\end{proposition}

This categorical structure parallels the category of field bundles over
$\mathcal{M}$: it provides the algebraic analogue of geometric functoriality.

% ============================================================
\section{The Free Typed Hypergraph Monad}
% ============================================================

We now construct the free hypergraph monad $\mathbb{F}$ on the category
$\mathbf{Type}$.

Given a collection of typed atoms $\{a_i\}$, the monad $\mathbb{F}$ freely
generates the smallest typed hypergraph containing them.

\begin{definition}[Free Typed Hypergraph Monad]
The endofunctor
\[
\mathbb{F}: \mathbf{Type} \to \mathbf{TypHyper}
\]
assigns to each type $\tau$ the free hypergraph generated by all atoms of type
$\tau$ and their admissible hyperedges.  Its unit $\eta$ inserts atoms, and its
multiplication $\mu$ performs admissible rewiring.
\end{definition}

\begin{proposition}
Algebras over $\mathbb{F}$ correspond to typed hypergraphs equipped with a
canonical content-atom insertion operation.
\end{proposition}

Thus the monadic structure encodes the generative rules of Polyxan media.

\begin{theorem}[Initial Algebra]
Content atoms form the initial $\mathbb{F}$-algebra.  Every typed hypergraph is
obtained by iterated insertion of content atoms followed by closure under
admissible hyperedges.
\end{theorem}

This parallels how fields on $\mathcal{M}$ arise from free sections of the jet
bundle (Chapter~\ref{ch:jet-bundles}).

% ============================================================
\section{The Semantic Tag Sheaf}
% ============================================================

For each typed hypergraph $\mathcal{G}$, we build a sheaf over its incidence
complex $\mathsf{Inc}(\mathcal{G})$.

\begin{definition}[Semantic Tag Sheaf]
The semantic tag sheaf $\mathscr{T}_{\mathcal{G}}$ assigns:
\[
\mathscr{T}_{\mathcal{G}}(U)
=
\{\text{tag assignments on } U \subseteq \mathsf{Inc}(\mathcal{G})\}
\]
with restriction given by tag refinement.
\end{definition}

\begin{proposition}
$\mathscr{T}_{\mathcal{G}}$ is a flabby sheaf, hence acyclic.
\end{proposition}

This sheaf captures local semantic consistency and forms the discrete analogue
of the entropy field $S$ on $\mathcal{M}$.

% ============================================================
\section{The Polyxan Curvature Functional}
% ============================================================

We now introduce the first Polyxan curvature functional, which measures the
local semantic deficit and tag entropy.

Let $\delta(e)$ denote the \emph{deficit} of a hyperedge $e$ (difference between
expected and actual typed arity), and let
\[
H_{\mathrm{tag}}(v) = -\sum_{\theta\in \mathbf{Tag}} p_\theta(v)\log p_\theta(v)
\]
be the tag entropy at vertex $v$.

\begin{definition}[Primitive Polyxan Curvature]
The primitive Polyxan curvature functional is
\[
\mathcal{H}_0[\mathcal{G}]
=
\sum_{e\in E} \delta(e)
\;+\;
\sum_{v\in V} H_{\mathrm{tag}}(v).
\]
\end{definition}

This quantity plays the rôle of the scalar curvature $R$ in Part~I.
It will serve as the discrete Lyapunov functional driving Polyxan reduction.

% ============================================================
\section{Media Quines as Fixed Points of the Free-Algebra Endofunctor}
% ============================================================

We now state the discrete analogue of RSVP vacuum states.

\begin{definition}[Media Quine]
A media quine is a typed hypergraph $\mathcal{G}$ satisfying
\[
\mathbb{F}(\mathcal{G}) \cong \mathcal{G},
\]
i.e.~a fixed point of the free typed hypergraph monad.
\end{definition}

\begin{theorem}[Quine Fixed-Point Characterization]
A typed hypergraph $\mathcal{G}$ is a media quine if and only if it is a
critical point of the curvature functional $\mathcal{H}_0$ and is stable under
content-atom insertion.
\end{theorem}

This parallels the characterization of RSVP vacua as critical points of the
semantic action.

% ============================================================
\section{Summary}
% ============================================================

In this chapter we:
\begin{itemize}
\item defined content atoms as the primitive combinatorial units of Polyxan media,
\item constructed typed hypergraphs and the category $\mathbf{TypHyper}$,
\item introduced the free typed hypergraph monad and proved content atoms
  form the initial algebra,
\item defined the semantic tag sheaf over the incidence complex,
\item introduced the primitive Polyxan curvature functional $\mathcal{H}_0$,
\item characterized media quines as fixed points of the free-algebra
  endofunctor and as curvature-critical configurations.
\end{itemize}

Typed hypergraphs now form a complete algebraic foundation for the Polyxan
hyperstructure.  The next chapter deepens this structure through spans,
cospans, and the rewriting calculus.
